\documentclass[aspectratio=169]{ctexbeamer}
\usepackage[T1]{fontenc}
\usepackage{latexsym,amsmath,xcolor,multicol,booktabs,calligra}
\usepackage{graphicx,listings,stackengine}
\usepackage{hyperref}

\usepackage{PekingU}

% 去掉每节/子节前面的自动目录页
\AtBeginSection[]{}
\AtBeginSubsection[]{}

% ===== 每节课改这里 =====
\author{Cap1taL}
\title{比赛题讲解}
\subtitle{?!比赛?!}
\date{2026年7月}
% ========================

% 代码高亮设置
\definecolor{deepblue}{rgb}{0,0,0.5}
\definecolor{deepred}{rgb}{0.6,0,0}
\definecolor{deepgreen}{rgb}{0,0.5,0}
\definecolor{halfgray}{gray}{0.55}

\lstset{
    basicstyle=\ttfamily\small,
    keywordstyle=\bfseries\color{deepblue},
    emphstyle=\ttfamily\color{deepred},
    stringstyle=\color{deepgreen},
    numbers=left,
    numberstyle=\small\color{halfgray},
    rulesepcolor=\color{red!20!green!20!blue!20},
    frame=shadowbox,
}

\begin{document}

\kaishu

% 封面
\begin{frame}
    \titlepage
\end{frame}

% 目录
% \begin{frame}{目录}
%     \tableofcontents[sectionstyle=show,subsectionstyle=show/shaded/hide,subsubsectionstyle=show/shaded/hide]
% \end{frame}

% ===== 从这里开始写你的内容 =====

\section{模拟赛}

\subsection{GCJ 2010 Theme Park}
\begin{frame}{GCJ 2010 Theme Park}
    过山车每次最多载 $k$ 人，运行 $R$ 次。$N$ 个小组排队，每组 $g_i$ 人。每次按顺序让尽可能多小组上车，满员或无法容纳下一组即出发。上车的小组按原顺序回到队尾。求总收入（每人 $1$ 欧元）。
    
    $R\le 10^8,\ N\le 1000, T\le 50$
\end{frame}

\begin{frame}{Solution}
    % 整行居中布局
    \centering
    % 左侧图片区域 占 0.48 文本宽
    \begin{minipage}{0.40\textwidth}
    \includegraphics[width=\linewidth]{pic/1.png}
    \end{minipage}
    \hfill % 中间留白分隔
    % 右侧文字区域
    \begin{minipage}{0.55\textwidth}
        \begin{itemize}
            \item 暴力模拟，一旦出现循环（循环定义为一个队列局面出现过）则无需继续计算
            \item 注意到一共有 $n$ 种局面
            \item 复杂度 $O(Tn^2)$
            \item 双指针预处理所有局面的上车情况，可以随便做到线性，$O(Tn)$
            \item 本质上是一个类似基环树森林的东西
            \item 这个题用时不应该超过 $20\min$，哪怕放在 NOIP 也是很简单的签到了
        \end{itemize}
    \end{minipage}
\end{frame}





\subsection{洛谷 P16928 聚魔石}
\begin{frame}{洛谷 P16928 聚魔石}
    $n$ 个聚魔石排成一排，编号 $1\sim n$，每个有魔力存储量 $p_i$（$1\sim n$ 排列）和存储容易度 $w_i$。
    
    初始 $s$ 充能，随后从已充能区间 $[l,r]$ 向相邻未充能位置扩展：若 $l>1$ 且 $r<n$，则以 $w_{l-1}/(w_{l-1}+w_{r+1})$ 概率向左、$w_{r+1}/(w_{l-1}+w_{r+1})$ 概率向右。
    
    若新充能石的 $p$ 严格大于此前所有已充能石的 $p$，则发生一次共鸣。初始 $s$ 也算一次共鸣。求共鸣期望次数，对 $998244353$ 取模。
    
    $n\le 2000$
\end{frame}

\begin{frame}{Solution}
    % 整行居中布局
    \centering
    % 左侧图片区域 占 0.48 文本宽
    \begin{minipage}{0.40\textwidth}
    
\includegraphics[width=\linewidth]{pic/2.jpg}
    \end{minipage}
    \hfill % 中间留白分隔
    % 右侧文字区域
    \begin{minipage}{0.55\textwidth}
        \begin{itemize}
            \item 区间不断扩展，自然想到区间 DP
            \item 设 $dp_{l,r}$ 表示现在 $[l,r]$ 已注魔，到住满还能期望共鸣几次
            \item 转移直接考虑左右即可 
            \item 这个题在熟悉期望的情况下是可以秒出做法的，比较模板
        \end{itemize}
    \end{minipage}
\end{frame}








\subsection{洛谷 P3538 [POI 2012] OKR-A Horrible Poem}

\begin{frame}{洛谷 P3538 [POI 2012] OKR-A Horrible Poem}

给定长度为 $n$ 的字符串，$q$ 次询问子串 $[l,r]$ 的最短完整周期长度。

\vspace{0.5em}
$1\le n\le 500\,000,\quad 1\le q\le 2\,000\,000$

\end{frame}



\begin{frame}{Solution}
    % 整行居中布局
    \centering
    % 左侧图片区域 占 0.48 文本宽
    \begin{minipage}{0.40\textwidth}
    \includegraphics[width=\linewidth]{pic/1.png}
    \end{minipage}
    \hfill % 中间留白分隔
    % 右侧文字区域
    \begin{minipage}{0.55\textwidth}
        \begin{itemize}
            \item 先考虑怎么判断一个长度是否是循环节
            \item 哈希比对 $[l,r-d]$ 和 $[l+d,r]$
            \item 现在问题在与要判断的长度太多。发现只有长度的约数才有可能是循环节，暴力枚举约数即可
            \item 复杂度 $O(qd)$，$d$ 是 $1~n$ 最大约数数量，实测是 $ 100 $。
        \end{itemize}
    \end{minipage}
\end{frame}




\subsection{洛谷 P4219 [BJOI2014] 大融合}

\begin{frame}{洛谷 P4219 [BJOI2014] 大融合}

    $N$ 个孤立星球，边按顺序添加，每次连边保证两端不连通。动态询问某条边当前的负载——即所在的树上经过该边的简单路径数量。

    \vspace{0.5em}
    $1\le N,Q\le 10^5$

\end{frame}



\begin{frame}{Solution}
    % 整行居中布局
    \centering
    % 左侧图片区域 占 0.48 文本宽
    \begin{minipage}{0.40\textwidth}
    \includegraphics[width=\linewidth]{pic/1.png}
    \end{minipage}
    \hfill % 中间留白分隔
    % 右侧文字区域
    \begin{minipage}{0.55\textwidth}
        \begin{itemize}
            \item 做法很多
            \item 首先发现，题目要计算的东西，等价于断开这条边后 $(u,v)$ 两点联通块大小的乘积
            \item 因此我们在每次查询 $(u,v)$ 之前，加一个移除操作，在查询之后再加回来
            \item 现在题面变成，动态加边删边，查询操作变为查询点所在的联通块大小
            \item 一眼线段树分治，之后变成加边和撤销边，查询联通块大小，使用可撤销并查集维护，复杂度 $O(q\log q)$
            \item 期望用时 $ 1 h $ 左右 
        \end{itemize}
    \end{minipage}
\end{frame}

\begin{frame}
    概率与期望？
    \begin{itemize}
        \item 条件概率？
        \item 期望的线性性？
        \item 全概率公式？
        \item 全期望公式？
        \item 二者的区别？
    \end{itemize}
\end{frame}

\subsection{Codeforces 280C Game on Tree}

\begin{frame}{Codeforces 280C Game on Tree}
      给定一棵 $n$ 个结点的有根树（根为 $1$）。每次等概率选择一个未被删除的结点，将其及其子树全部删去。求删除所有结点的期望操作次数。
      
      $1\le n\le 10^5$
\end{frame}
% ===== 结束页 =====

\begin{frame}
    \begin{center}
        {\Huge\calligra Thanks!}
    \end{center}
\end{frame}

\end{document}
